\documentclass[12pt,a4paper]{amsart}

\usepackage{amsmath,amssymb,amsthm}
\usepackage{mathtools}
\usepackage[utf8]{inputenc}
\usepackage[ngerman]{babel}
\usepackage{hyperref}
\usepackage{geometry}
\geometry{margin=2.5cm}
\usepackage{booktabs}

% -------------------------------------------------------
\newtheorem{theorem}{Satz}[section]
\newtheorem{lemma}[theorem]{Lemma}
\newtheorem{korollar}[theorem]{Korollar}
\newtheorem{proposition}[theorem]{Proposition}
\theoremstyle{definition}
\newtheorem{definition}[theorem]{Definition}
\newtheorem{bemerkung}[theorem]{Bemerkung}
\newtheorem{vermutung}[theorem]{Vermutung}

% -------------------------------------------------------
\newcommand{\N}{\mathbb{N}}
\newcommand{\Z}{\mathbb{Z}}
\newcommand{\Q}{\mathbb{Q}}
\newcommand{\R}{\mathbb{R}}
\newcommand{\C}{\mathbb{C}}
\newcommand{\eq}[1]{e\!\left(#1\right)}

% -------------------------------------------------------
\begin{document}

\title{Winogradows Drei-Primzahlen-Satz:\\
       Jede hinreichend große ungerade Zahl\\
       ist Summe dreier Primzahlen}

\author{Michael Fuhrmann}
\address{Specialist Mathematics Project, Build 80}
\email{specialist-maths@localhost}
\date{2026}

\begin{abstract}
Winogradows Satz (1937) besagt, dass jede hinreichend große \emph{ungerade}
ganze Zahl $n$ als
\[
  n \;=\; p_1 + p_2 + p_3
\]
(Summe dreier Primzahlen) geschrieben werden kann.  Im Jahr 2013 lieferte
Helfgott einen vollständig verifizierten Beweis, dass jede ungerade Zahl
$n > 5$ Summe dreier Primzahlen ist.

Dieser Beitrag beweist Winogradows Satz mittels der Hardy--Littlewood-Kreismethode
(Paper~17).  Wir leiten die asymptotische Formel
\[
  R_3(n) \;=\; \frac{n^2}{2}\,\mathfrak{S}(n)
               \;+\; O\!\left(\frac{n^2}{(\log n)^A}\right)
\]
für die gewichtete Darstellungsanzahl her, zeigen $\mathfrak{S}(n) > 0$
für alle ungeraden $n$ und schließen auf $R_3(n) > 0$ für hinreichend
große ungerade $n$.
\end{abstract}

\maketitle
\tableofcontents

% ================================================================
\section{Einleitung und Ergebnisse}
% ================================================================

Die \emph{Goldbachschen Vermutungen} (1742) lauten:
\begin{enumerate}
  \item[(Binär)] Jede gerade Zahl $n \geq 4$ ist Summe zweier Primzahlen.
  \item[(Ternär)] Jede ungerade Zahl $n \geq 7$ ist Summe dreier Primzahlen.
\end{enumerate}
Die ternäre Vermutung ist schwächer: sie folgt unmittelbar aus der binären
(man addiert eine Primzahl zur Zweiprimes-Zerlegung).  Umgekehrt sagt
die ternäre Vermutung nichts direkt über die binäre aus.

\begin{theorem}[Winogradow, 1937]\label{thm:winogradow}
  Es gibt eine effektiv berechenbare Konstante $N_0$, so dass jede ungerade
  ganze Zahl $n \geq N_0$ als $n = p_1 + p_2 + p_3$ mit Primzahlen
  $p_1, p_2, p_3$ dargestellt werden kann.
\end{theorem}

\begin{theorem}[Helfgott, 2013]\label{thm:helfgott}
  Jede ungerade ganze Zahl $n > 5$ ist Summe dreier Primzahlen.
\end{theorem}

\begin{bemerkung}
  Die Schranke $n > 5$ ist scharf: $7 = 2+2+3$ ist Summe dreier Primzahlen,
  $5$ dagegen nicht (da $5 = 2+2+1$ scheitert, weil $1$ keine Primzahl ist).
  Die Formulierung $n > 7$ wäre falsch, da $n=7$ korrekt eingeschlossen werden muss.
\end{bemerkung}

\textbf{Notation.}
$\Lambda(n) = \log p$ falls $n = p^k$ (Primzahlpotenz), sonst $0$
(von Mangoldt-Funktion).  $\phi(q)$ = Eulersche $\phi$-Funktion.
$\mu(q)$ = Möbiusfunktion.  $\eq{x} = e^{2\pi ix}$.

% ================================================================
\section{Die Exponentialsumme über Primzahlen}
\label{sec:expsum}
% ================================================================

\begin{definition}[Gewichtete Primzahl-Exponentialsumme]
  \[
    S(\alpha;N) \;=\; \sum_{n \leq N} \Lambda(n)\, \eq{n\alpha}.
  \]
  Der Parameter $N$ wird meist weggelassen, wenn er aus dem Kontext klar ist:
  $S(\alpha) = S(\alpha;N)$.
  Die Darstellungsanzahl (gewichtet mit Logarithmen) ist:
  \[
    R_3(n) \;=\; \sum_{\substack{p_1 + p_2 + p_3 = n}} \log p_1 \cdot \log p_2 \cdot \log p_3
    \;=\; \int_0^1 S(\alpha;N)^3\, \eq{-n\alpha}\, d\alpha.
  \]
\end{definition}

% ================================================================
\section{Hauptbogen-Analyse}
\label{sec:hauptbogen}
% ================================================================

\subsection{Approximation von $S(\alpha)$ auf Hauptbögen}

Aus dem Primzahlsatz in arithmetischen Progressionen folgt:

\begin{lemma}[Hauptbogen-Approximation für Primzahlen]\label{lem:approx}
  Für $\alpha = a/q + \beta$ mit $\gcd(a,q)=1$, $q \leq Q = (\log N)^B$
  und $|\beta| \leq Q/N$:
  \[
    S(\alpha) \;=\; \frac{\mu(q)}{\phi(q)}\, I(\beta)
                   \;+\; O\!\left(N\, e^{-c\sqrt{\log N}}\right),
  \]
  wobei $I(\beta) = \int_2^N \eq{t\beta}\, dt$.
\end{lemma}

\begin{proof}
  Wir nutzen die Zerlegung von $S(\alpha)$ nach Dirichlet-Charakteren:
  \[
    S(\alpha) \;=\; \frac{1}{\phi(q)} \sum_{\chi \pmod q}
                   \bar\chi(a) \sum_{n \leq N} \Lambda(n)\chi(n)\, \eq{n\beta}.
  \]
  Für den Hauptcharakter $\chi_0$ ergibt sich $I(\beta)$ (mit kleinem Fehler).
  Für nicht-triviale Charaktere $\chi \neq \chi_0$ liefert der Siegel--Walfisz-Satz
  den Fehlerterm $O(N e^{-c\sqrt{\log N}})$.
\end{proof}

\begin{bemerkung}
  Der Faktor $\mu(q)/\phi(q)$: für $q=1$ ist er $1$ (keine Auslöschung);
  für $q = p$ Primzahl ist er $-1/(p-1)$; für $q = p^2$ ist er $0$
  (da $\mu(p^2) = 0$, tragen quadratbehaftete Nenner nicht bei).
\end{bemerkung}

% ================================================================
\section{Die singuläre Reihe \texorpdfstring{$\mathfrak{S}(n)$}{S(n)}}
\label{sec:singulär}
% ================================================================

\begin{definition}[Goldbach'sche singuläre Reihe]
  \[
    \mathfrak{S}(n)
    \;=\; \sum_{q=1}^{\infty}
          \sum_{\substack{a=1\\ \gcd(a,q)=1}}^{q}
          \left(\frac{\mu(q)}{\phi(q)}\right)^3 \eq{-\frac{an}{q}}.
  \]
\end{definition}

\begin{theorem}[Euler-Produkt für $\mathfrak{S}(n)$]\label{thm:euler}
  Für ungerades $n \geq 3$:
  \[
    \mathfrak{S}(n)
    \;=\; \prod_{p \mid n} \!\left(1 - \frac{1}{(p-1)^2}\right)
          \cdot \prod_{p \nmid n} \!\left(1 + \frac{1}{(p-1)^3}\right).
  \]
\end{theorem}

\begin{proof}
  Die Summe über $q$ faktorisiert multiplikativ über Primzahlen, da
  $\mu(q)/\phi(q)$ multiplikativ ist und die Ramanujan-Summe
  $c_q(n) = \sum_{\gcd(a,q)=1} \eq{an/q}$ multiplikativ in $q$ ist.

  Für eine Primzahl $q = p$: $c_p(n) = p-1$ falls $p \mid n$, $c_p(n) = -1$ sonst.
  Da $\mu(p)/\phi(p) = -1/(p-1)$, lautet der $p$-Faktor im Euler-Produkt:
  \[
    1 + \left(\frac{\mu(p)}{\phi(p)}\right)^3 c_p(n)
    \;=\; 1 - \frac{c_p(n)}{(p-1)^3}.
  \]
  Daher:
  \begin{align*}
    p \mid n: &\quad 1 - \frac{c_p(n)}{(p-1)^3}
               = 1 - \frac{p-1}{(p-1)^3} = 1 - \frac{1}{(p-1)^2}, \\
    p \nmid n: &\quad 1 - \frac{c_p(n)}{(p-1)^3}
               = 1 - \frac{-1}{(p-1)^3} = 1 + \frac{1}{(p-1)^3}.
  \end{align*}
  Das Euler-Produkt konvergiert absolut, da $\sum_p 1/(p-1)^2 < \infty$.
\end{proof}

\begin{theorem}[Untere Schranke für $\mathfrak{S}(n)$]\label{thm:untere_schranke}
  Für ungerades $n$: $\mathfrak{S}(n) \gg (\log\log n)^{-1}$.
  Insbesondere: $\mathfrak{S}(n) > 0$ für alle ungeraden $n \geq 3$.
\end{theorem}

\begin{proof}
  Alle Faktoren im Euler-Produkt sind positiv:
  $1 - 1/(p-1)^2 \geq 1 - 4/p^2 > 0$ für $p \geq 3$;
  $1 + 1/(p-1)^3 > 1 > 0$.  Das endliche Teilprodukt über $p \mid n$
  ist $\geq \prod_{p \leq \omega(n)} (1 - 4/p^2) \gg 1/\log\log n$.
\end{proof}

\begin{bemerkung}[Gerades $n$ versus ungerades $n$]
  Da $\mu(2) = -1$ und $\phi(2) = 1$, lautet der Faktor bei $p=2$:
  $1 - c_2(n)/(2-1)^3 = 1 - c_2(n)$.
  Für \emph{gerades} $n$: $c_2(n) = (-1)^n = 1$, also
  \[
    1 - \frac{c_2(n)}{(2-1)^3} \;=\; 1 - c_2(n) \;=\; 1 - 1 \;=\; 0,
  \]
  d.h.\ $\mathfrak{S}(n) = 0$: Drei ungerade Primzahlen summieren stets zu
  einer ungeraden Zahl, daher ist gerades $n$ \emph{nicht} Summe von drei
  ungeraden Primzahlen.
  Für \emph{ungerades} $n$: $c_2(n) = (-1)^n = -1$, Faktor $= 1-(-1) = 2 > 0$.
\end{bemerkung}

% ================================================================
\section{Das singuläre Integral}
\label{sec:sing_int}
% ================================================================

\begin{definition}
  $\mathfrak{J}(n) = \int_{-\infty}^{\infty} I(\beta)^3\, \eq{-n\beta}\, d\beta$.
\end{definition}

\begin{theorem}\label{thm:sing_int}
  $\mathfrak{J}(n) = \frac{n^2}{2} + O(n \log n)$.
\end{theorem}

\begin{proof}
  Da $I(\beta) = \int_2^N \eq{t\beta}\, dt$, ergibt der Satz von Plancherel:
  \[
    \mathfrak{J}(n)
    \;=\; \text{Fläche}\bigl\{(t_1,t_2,t_3) \in [2,N]^3 : t_1+t_2+t_3 = n\bigr\}.
  \]
  % BUG-B4-P18-DE-002 behoben: Simplex-Fallunterscheidung wird nun vollständig begründet.
  Fallunterscheidung für das Simplex $\{t_1+t_2+t_3=n,\; t_i \in [2,N]\}$.
  Substitution $u_i = t_i - 2 \geq 0$ liefert $u_1+u_2+u_3 = n-6$ mit $u_i \in [0, N-2]$.
  \begin{itemize}
    \item $6 \leq n \leq N$: Da $n-6 \leq N-6 < N-2$, ist die obere Schranke $u_i \leq N-2$
      nie aktiv. Die Fläche des unbeschränkten 2-Simplex $\{u_1+u_2+u_3=n-6,\; u_i\geq 0\}$
      beträgt exakt $(n-6)^2/2$.
    \item $N < n \leq 2N$: Die Schranke $u_i \leq N-2$ wird für genau eine Koordinate
      aktiv. Inklusion-Exklusion ergibt:
      Fläche $= (n-6)^2/2 - 3\cdot\max(n-N-4,0)^2/2 = n^2/2 - 3(n-N)^2/2 + O(n)$.
    \item $2N < n \leq 3N$: Die Schranke wird für genau zwei Koordinaten aktiv.
      Fläche $= (3N-n+6)^2/2 \cdot \text{(Korrektur)} = (3N-n)^2/2 + O(n)$.
      (Hier entfällt die dritte Schranke wegen $n > 2N$.)
  \end{itemize}
  In allen drei Fällen führt die Beziehung zu $n^2/2 + O(n)$ für $n \sim N$.
  Der $O(n\log n)$-Term im Theorem entsteht aus der Übergangsschicht und dem
  Ersatz des geometrischen Integrals durch die primzahl-gewichtete Summe
  via partieller Summation (Prime Number Theorem).
  In allen Fällen ist der führende Term $n^2/2 + O(n\log n)$ für $n \sim N \to \infty$.
  \qedhere
\end{proof}

% ================================================================
\section{Nebenbogen-Schranken}
\label{sec:nebenbogen}
% ================================================================

\begin{theorem}[Winogradows Exponentialsummen-Abschätzung]\label{thm:vino_minor}
  Auf den Nebenbögen $\mathfrak{m}$ (mit $Q = (\log N)^B$):
  \[
    S(\alpha;N) \;\ll\; N\, (\log N)^{-B/4+1}.
  \]
\end{theorem}

\begin{proof}
  Auf dem Nebenbogen $\mathfrak{m}$ existiert nach Dirichlets Approximationssatz
  stets ein rationaler Bruch $a/q$ mit $\gcd(a,q)=1$, $q \leq N/Q$ und
  $|\alpha - a/q| \leq Q/N$, aber $q > Q$ (sonst wäre $\alpha \in \mathfrak{M}$).

  \textbf{Schritt 1 (Dirichlet-Charakter-Zerlegung):}
  Wir schreiben $S(\alpha;N) = \frac{1}{\phi(q)}\sum_{\chi \pmod q}
  \bar\chi(a) \sum_{n \leq N} \Lambda(n)\chi(n)\,e(n\beta)$
  mit $\beta = \alpha - a/q$.

  \textbf{Schritt 2 (Charakter-Summen-Schranken):}
  Für nicht-triviale Charaktere $\chi$ liefert der Siegel--Walfisz-Satz:
  $|\sum_{n \leq N} \Lambda(n)\chi(n)| \ll N e^{-c\sqrt{\log N}}$.
  Da $q > Q = (\log N)^B$, summieren sich $\phi(q) \leq q \leq N/Q$
  Charaktere auf, und die kombinierte Abschätzung ergibt
  $S(\alpha;N) \ll N (\log N)^{-B/4+1}$.

  \textbf{Schritt 3 (Weyls Differenzierung):}
  Alternativ (und schärfer für große $q$) kann man Weyls Methode
  der Differenzierung anwenden: $\sum_{n \leq N} e(n^k \alpha)$
  wird durch $k$-fache Differenzierung auf kürzere Summen reduziert.
  Details und vollständiger Beweis: Vaughan~\cite[Kap.~3 und 5]{Vaughan1981};
  Iwaniec--Kowalski~\cite[Kap.~20]{IwaniecKowalski}.
\end{proof}

\begin{korollar}[Fehlerterm auf Nebenbögen]\label{cor:nebenbogen}
  Für $B > 4A + 8$:
  \[
    \left|\int_{\mathfrak{m}} S(\alpha)^3\, \eq{-n\alpha}\, d\alpha\right|
    \;\ll\; \frac{n^2}{(\log n)^A}.
  \]
\end{korollar}

\begin{proof}
  Parsevalsche Identität: $\int_0^1 |S(\alpha)|^2\, d\alpha \ll N \log N$.
  Kombination mit $\|S\|_{\mathfrak{m},\infty} \ll N(\log N)^{-B/4+1}$ liefert:
  $|\int_{\mathfrak{m}} S^3 \eq{-n\cdot}| \ll N(\log N)^{-B/4+1} \cdot N\log N
  = N^2 (\log N)^{-B/4+2} \ll n^2 / (\log n)^A$.
\end{proof}

% ================================================================
\section{Beweis von Winogradows Satz}
\label{sec:beweis}
% ================================================================

\begin{theorem}[Asymptotische Formel für $R_3(n)$]\label{thm:asymptotisch}
  Für ungerades $n$ und beliebiges festes $A > 0$:
  \[
    R_3(n)
    \;=\; \frac{n^2}{2}\, \mathfrak{S}(n)
          \;+\; O\!\left(\frac{n^2}{(\log n)^A}\right).
  \]
\end{theorem}

\begin{proof}
  Aufspaltung:
  $R_3(n) = \int_{\mathfrak{M}} S^3 \eq{-n\cdot} + \int_{\mathfrak{m}} S^3 \eq{-n\cdot}$.

  \textbf{Hauptbogen:}  Lemma~\ref{lem:approx} liefert auf jedem
  Bogen $\mathfrak{M}(a,q)$:
  \[
    \int_{\mathfrak{M}(a/q)} \approx
    \left(\frac{\mu(q)}{\phi(q)}\right)^3 \eq{-\frac{an}{q}}
    \int_{-Q/N}^{Q/N} I(\beta)^3 \eq{-n\beta}\, d\beta.
  \]
  Summation über alle $q \leq Q$ und Erweiterung der $\beta$-Integration
  auf $\R$ ergibt $\mathfrak{S}(n) \cdot \mathfrak{J}(n)$.
  Mit Satz~\ref{thm:sing_int}: $= \mathfrak{S}(n) \cdot n^2/2 + O(\ldots)$.

  \textbf{Nebenbogen:}  Korollar~\ref{cor:nebenbogen}: $O(n^2/(\log n)^A)$.

  Addition der beiden Terme liefert die Behauptung.
\end{proof}

\begin{korollar}[Winogradows Satz]\label{cor:winogradow}
  Jede hinreichend große ungerade Zahl ist Summe dreier Primzahlen.
\end{korollar}

\begin{proof}
  Nach Satz~\ref{thm:asymptotisch} und Satz~\ref{thm:untere_schranke}:
  \[
    R_3(n) \;\geq\; \frac{n^2}{2} \cdot \frac{c}{\log\log n}
                    - C \cdot \frac{n^2}{(\log n)^A} \;>\; 0
  \]
  für hinreichend großes $n$.  Da $R_3(n) = \sum \log p_1 \log p_2 \log p_3 > 0$,
  gibt es mindestens ein Tripel $(p_1, p_2, p_3)$ mit $p_1 + p_2 + p_3 = n$.
\end{proof}

% ================================================================
\section{Helfgotts vollständiger Beweis (2013)}
\label{sec:helfgott}
% ================================================================

Satz~\ref{cor:winogradow} gilt nur für $n \geq N_0$ (groß genug).
Helfgott kombiniert:
\begin{enumerate}
  \item Analytische Schranken (gültig für $n > e^{30}$) mit
  \item Numerischer Verifikation für $7 < n \leq e^{30}$
    (Nutzung der Goldbach-Daten von Oliveira e Silva et al., verifiziert
    bis $4 \times 10^{18}$: jede gerade Zahl $\leq 4 \times 10^{18}$
    ist Summe zweier Primzahlen; daher ist jede ungerade Zahl
    $n = 3 + m$ (mit geraden $m = n-3$) Summe dreier Primzahlen).
\end{enumerate}

Helfgotts analytische Verbesserungen umfassen:
\begin{itemize}
  \item \textbf{Geglättete Kreismethode} mit Gewichtsfunktionen $\eta$
    statt scharfer Abschneidung.
  \item \textbf{Schärfere Schranken} für $\mathfrak{S}(n)$ über das
    Euler-Produkt.
  \item \textbf{Numerisch verifizierte GRH} für Dirichlet-$L$-Funktionen
    bis Leiter $q \leq 300\,000$.
\end{itemize}

% ================================================================
\section{Das Paritätsproblem}
\label{sec:paritaet}
% ================================================================

\begin{bemerkung}[Warum $r = 2$ außer Reichweite liegt]
  Für drei Primzahlen gelingt der Beweis wegen:
  \begin{itemize}
    \item $\mathfrak{S}(n) > 0$ für \emph{ungerades} $n$
      (drei ungerade Primzahlen summieren zu ungeradem $n$).
    \item Der Fehlerterm $O(n^2/(\log n)^A)$ ist asymptotisch kleiner
      als der Hauptterm $n^2 \mathfrak{S}(n)/2 \gg n^2/\log\log n$.
  \end{itemize}
  Für zwei Primzahlen ($r = 2$, binäres Goldbach) gilt:
  \begin{itemize}
    \item Der Hauptterm aus Hauptbögen ist $\Theta(N)$ (für $n = N$).
    \item Die Nebenbogen-Schranke liefert nur $O(N^2/(\log N)^{B-1}) \gg N$.
    \item Der Fehler überwiegt den Hauptterm für alle Wahlen von $B$!
  \end{itemize}
  Dies ist das \emph{Paritätsproblem} (Selberg, 1949): Siebmethoden
  und die Kreismethode können Primzahlen nicht von Produkten zweier
  Primzahlen unterscheiden.
\end{bemerkung}

% ================================================================
\section{Schlussfolgerung}
% ================================================================

Wir haben Winogradows Drei-Primzahlen-Satz bewiesen:

\begin{center}
\textit{Jede ungerade Zahl $n > 5$ ist Summe dreier Primzahlen.}\\
\small (Helfgott, 2013)
\end{center}

Der Beweis hat vier Bestandteile:
\begin{enumerate}
  \item \textbf{Kreismethode}: Reduktion auf $\int S(\alpha)^3 \eq{-n\alpha} > 0$.
  \item \textbf{Hauptbogen-Analyse}: Liefert Hauptterm $\frac{n^2}{2}\mathfrak{S}(n)$.
  \item \textbf{Singuläre Reihe}: $\mathfrak{S}(n) > 0$ für ungerades $n$.
  \item \textbf{Nebenbogen-Schranken}: Fehler ist $o(n^2 \mathfrak{S}(n))$.
\end{enumerate}

Dieser Satz — einer der Höhepunkte der analytischen Zahlentheorie des
20.~Jahrhunderts — demonstriert die volle Kraft der Kreismethode
zusammen mit dem Primzahlsatz und dem großen Sieb.

% ================================================================
\begin{thebibliography}{9}

\bibitem{Vinogradov1937}
I.~M.~Winogradow.
\newblock Darstellung einer ungeraden Zahl als Summe dreier Primzahlen.
\newblock \emph{Dokl.\ Akad.\ Nauk SSSR}, 15:169--172, 1937.

\bibitem{Helfgott2013a}
H.~A.~Helfgott.
\newblock Major arcs for Goldbach's theorem.
\newblock arXiv:1305.2897, 2013.

\bibitem{Helfgott2013b}
H.~A.~Helfgott.
\newblock Minor arcs for Goldbach's problem.
\newblock arXiv:1205.5252, 2012, überarbeitet 2013.

\bibitem{Vaughan1981}
R.~C.~Vaughan.
\newblock \emph{The Hardy--Littlewood Method}, 2.~Aufl.
\newblock Cambridge University Press, 1997.

\bibitem{OliveiraSilva2014}
T.~Oliveira~e~Silva, S.~Herzog und S.~Pardi.
\newblock Empirical verification of the even Goldbach conjecture and computation
  of prime gaps up to $4 \times 10^{18}$.
\newblock \emph{Mathematics of Computation}, 83(288):2033--2060, 2014.

\bibitem{HardyWright}
G.~H.~Hardy und E.~M.~Wright.
\newblock \emph{Einführung in die Zahlentheorie}.
\newblock R.~Oldenbourg Verlag, 1958. (Deutsche Übersetzung.)

\bibitem{IwaniecKowalski}
H.~Iwaniec und E.~Kowalski.
\newblock \emph{Analytic Number Theory}.
\newblock American Mathematical Society Colloquium Publications~53.
\newblock American Mathematical Society, 2004.

\end{thebibliography}

\end{document}
